\begin{theorem}[Triangle Inequality]
Let $(V,\langle\cdot,\cdot\rangle)$ be an inner‑product space and let $x,y\in V$. Then
\[
\|x+y\|\le \|x\|+\|y\|.
\]
\end{theorem}

\begin{proof}
\textbf{1. Norm expressed through the inner product.} By definition of the norm induced by the inner product,
\[
\|v\|^{2}= \langle v,v\rangle \qquad\text{for every }v\in V .
\tag{1}
\]

\textbf{2. Expand $\|x+y\|^{2}$.} Using bilinearity in the first argument and conjugate‑linearity in the second argument we obtain
\begin{align*}
\|x+y\|^{2}
&=\langle x+y,x+y\rangle \\[2mm]
&=\langle x,x\rangle +\langle x,y\rangle +\langle y,x\rangle +\langle y,y\rangle .
\end{align*}
\tag{2}
Since $\langle y,x\rangle =\overline{\langle x,y\rangle}$, the sum of the mixed terms equals twice the real part of $\langle x,y\rangle$:
\[
\langle x,y\rangle +\langle y,x\rangle
= \langle x,y\rangle +\overline{\langle x,y\rangle}
= 2\,\operatorname{Re}\langle x,y\rangle .
\tag{3}
\]
Substituting (3) into (2) and using (1) for $\|x\|^{2}$ and $\|y\|^{2}$ gives
\[
\|x+y\|^{2}= \|x\|^{2}+2\,\operatorname{Re}\langle x,y\rangle+\|y\|^{2}.
\tag{4}
\]

\textbf{3. Bound the real part of the inner product.} The Cauchy–Schwarz inequality states
\[
|\langle x,y\rangle|\le \|x\|\,\|y\|.
\tag{5}
\]
Because the real part of any complex number does not exceed its modulus,
\[
\operatorname{Re}\langle x,y\rangle\le |\langle x,y\rangle|.
\tag{6}
\]
Combining (5) and (6) we obtain
\[
\operatorname{Re}\langle x,y\rangle\le \|x\|\,\|y\|.
\tag{7}
\]

\textbf{4. Insert the bound into (4).} From (4) and (7),
\[
\|x+y\|^{2}\le \|x\|^{2}+2\,\|x\|\,\|y\|+\|y\|^{2}
            = (\|x\|+\|y\|)^{2}.
\tag{8}
\]

\textbf{5. Take square roots.} Norms are non‑negative, i.e. $\|z\|\ge0$ for all $z\in V$. Hence both sides of (8) are non‑negative, and the monotonicity of the square‑root function on $[0,\infty)$ yields
\[
\|x+y\|\le \|x\|+\|y\|.
\tag{9}
\]

\medskip
\noindent\textbf{Equality case.} Equality in \eqref{9} occurs precisely when equality holds in both \eqref{5} and \eqref{6}.  
Equality in the Cauchy–Schwarz inequality holds iff $x$ and $y$ are linearly dependent, i.e. $y=\lambda x$ for some $\lambda\in\mathbb{R}$.  
Equality in \eqref{6} holds iff $\langle x,y\rangle$ is a non‑negative real number, which forces $\lambda\ge0$.  

Consequently, equality in the triangle inequality holds exactly when there exists a non‑negative real scalar $\lambda$ such that $y=\lambda x$ (including the trivial cases $x=0$ or $y=0$).
\end{proof}